% !TeX root = ../algebralineare.tex

\section{Matrici} \label{sec:matrices}
\subsection{Definizione di matrice} \label{subsec:matdefinition}
    Una matrice $A \in \mathbb{R}^{m, n}$ di dimensioni $m \times n$ a coefficienti reali è un oggetto matematico della forma:
    \[  A =
        \begin{bmatrix}
            a_{11} & a_{12} & ... & a_{1n} \\
            a_{21} & a_{22} & ... & a_{2n} \\
            \vdots & \ddots & \ddots & \vdots \\
            a_{m1} & a_{m2} & ... & a_{mn}
        \end{bmatrix}
    \]
    Dove $a_{ij} \in \mathbb{R}$ sono i coefficienti della matrice, $m$ è il numero delle righe ed $n$ il numero delle colonne.
    Se $n = m$ la matrice si dice quadrata e gli elementi $\{ a_{ii} \; | \; i \in \{1, ..., n\}\}$ sono chiamati ``diagonale'' della matrice.

    \vspace{10pt}
    Sia $\mathbb{R}^{m, n}$ l'insieme di tutte le generiche matrici $m \times n$, possiamo affermare che sussiste una biezione fra $\mathbb{R}^{m, 1} \leftrightarrow \mathbb{R}^{m}$ e fra $\mathbb{R}^{1, n} \leftrightarrow \mathbb{R}^{n}$.



\subsection{Prodotto scalare-matrice} \label{subsec:matscalarmul}
    Sia $A \in \mathbb{R}^{m, n}$ una matrice e sia $\lambda \in \mathbb{R}$ uno scalare, il prodotto scalare-matrice di $\lambda$ con $A$ è definito come:
    \[
    \lambda A =
    \begin{bmatrix}
        \lambda a_{11} & \lambda a_{12} & ... & \lambda a_{1n} \\
        \lambda a_{21} & \lambda a_{22} & ... & \lambda a_{2n} \\
        \vdots & \ddots & \ddots & \vdots \\
        \lambda a_{m1} & \lambda a_{m2} & ... & \lambda a_{mn}
    \end{bmatrix}
    \]

    \vspace{10pt}
    Proprietà:
    \begin{enumerate}[label=(\arabic*)]
        \item Associativa: $\forall \lambda, \mu \in \mathbb{R}, \; \forall A \in \mathbb{R}^{m, n}, \; (\lambda \mu) A = \lambda (\mu A)$
        \item Distributiva rispetto alla somma fra scalari: $\forall \lambda, \mu \in \mathbb{R}, \; \forall A \in \mathbb{R}^{m, n}, \; (\lambda + \mu)A = \lambda A + \mu A$
        \item Esistenza dell'elemento neutro: $\forall A \in \mathbb{R}^{m, n} \; \exists (1) \in \mathbb{R} \; | \; (1)A = A$
    \end{enumerate}



\subsection{Somma fra matrici} \label{subsec:matsum}
    Siano $A, B \in \mathbb{R}^{m, n}$ due matrici $m \times n$, la somma di $A$ con $B$ è definita come:
    \[
        \begin{aligned}
        A + B
        &=
        \begin{bmatrix}
            a_{11} & a_{12} & ... & a_{1n} \\
            a_{21} & a_{22} & ... & a_{2n} \\
            \vdots & \ddots & \ddots & \vdots \\
            a_{m1} & a_{m2} & ... & a_{mn}
        \end{bmatrix}
        +
        \begin{bmatrix}
            b_{11} & b_{12} & ... & b_{1n} \\
            b_{21} & b_{22} & ... & b_{2n} \\
            \vdots & \ddots & \ddots & \vdots \\
            b_{m1} & b_{m2} & ... & b_{mn}
        \end{bmatrix}
        \\
        &=
        \begin{bmatrix}
            a_{11} + b_{11} & a_{12} + b_{12} & ... & a_{1n} + b_{1n} \\
            a_{21} + b_{21} & a_{22} + b_{22} & ... & a_{2n} + b_{2n} \\
            \vdots & \ddots & \ddots & \vdots \\
            a_{m1} + b_{m1} & a_{m2} + b_{m2} & ... & a_{mn} + b_{mn}
        \end{bmatrix}
        \end{aligned}
    \]

    \vspace{10pt}
    Proprietà:
    \begin{enumerate}[label=(\arabic*)]
        \item Commutativa: $\forall A, B \in \mathbb{R}^{m, n}, \; A + B = B + A$
        \item Associativa: $\forall A, B, C \in \mathbb{R}^{m, n}, \; (A + B) + C = A + (C + B)$
        \item Distributiva del prodotto scalare-matrice rispetto alla somma matriciale: $\forall \lambda \in \mathbb{R}, \; \forall A, B \in \mathbb{R}^{m, n}, \; \lambda(A + B) = \lambda A + \lambda B$
        \item Esistenza dell'elemento neutro: $\exists 0_{m, n} \in \mathbb{R}^{m, n} \; | \; \forall A \in \mathbb{R}^{m, n}, \; A + 0_{m, n} = A $
        \item Esistenza dell'elemento opposto: $\forall A \in \mathbb{R}^{m, n}, \; \exists (-A) \in \mathbb{R}^{m, n} \; | \; A + (-A) = 0_{m, n}$
    \end{enumerate}



\subsection{Matrici trasposte} \label{subsec:mattranspose}
    Data una matrice $A \in \mathbb{R}^{m, n}$, la sua trasposta $A^{T} \in \mathbb{R}^{n, m}$ è la matrice nella forma:
    \[
        A =
        \begin{bmatrix}
            a_{11} & a_{12} & ... & a_{1n} \\
            a_{21} & a_{22} & ... & a_{2n} \\
            \vdots & \ddots & \ddots & \vdots \\
            a_{m1} & a_{m2} & ... & a_{mn}
        \end{bmatrix}
        \implies
        A^{T} =
        \begin{bmatrix}
            a_{11} & a_{21} & ... & a_{m1} \\
            a_{12} & a_{22} & ... & a_{m2} \\
            \vdots & \ddots & \ddots & \vdots \\
            a_{n1} & a_{n2} & ... & a_{nm}
        \end{bmatrix}
    \]


    Una matrice $A \in \mathbb{R}^{n, n}$ si dice simmetrica se $A = A^{T}$ oppure antisimmetrica se $A = -A^{T}$

    \vspace{10pt}
    \begin{theorem}
        Sia $V$ lo spazio vettoriale determinato da $(\mathbb{R}^{n, n}, \mathbb{R}, +, \cdot)$, allora
        \begin{enumerate}[label=(\Roman*)]
            \item $S = \{A \in \mathbb{R}^{n, n} \; | \; A = A^{T}\}$ e $W = \{A \in \mathbb{R}^{n, n} \; | \; A = -A^{T}\}$ sono sottospazi di $V$
            \item $S \oplus W = \mathbb{R}^{n, n}$
        \end{enumerate}

    \end{theorem}
    \begin{proof}[Dimostrazione]
        (I) Si consideri
        \[
            \forall A, B \in S, \; (A + B)^{T} = A^T + B^T = A + B \implies A + B \in S
        \]
        \[
            \forall A \in S, \; \forall \lambda \in \mathbb{R}, \; (\lambda A)^T = \lambda A^T = \lambda A \implies \lambda A \in S
        \]
        allora $S$ è sottospazio di $V$. Per $W$ invece
        \[
            \forall A, B \in W, \; (A + B)^{T} = A^T + B^T = -A + -B = -(A + B) \implies A + B \in W
        \]
        \[
            \forall A \in W, \; \forall \lambda \in \mathbb{R}, \; (\lambda A)^T = \lambda (A^T) = - \lambda A \implies \lambda A \in W
        \]
        allora $W$ è sottospazio di $V$.

        \vspace{10pt}
        (II) Se $A \in S \cap W$, si ha
        \[
            A = A^T = -A \implies A = 0_{n, n} \implies S \cap W = \{0_{n, n}\}
        \]
        quindi la somma $S \oplus W$ è diretta. Adesso, sia $M \in \mathbb{R}^{n, n}$. Se $A = \frac{M + M^{T}}{2} \in S$ e $B = \frac{M - M^{T}}{2} \in W$ abbiamo che $A + B = M \in S \oplus W$. Quindi $\mathbb{R}^{n, n} \subset S \oplus W$ e $S \oplus W \subset \mathbb{R}^{n, n} \implies S \oplus W = \mathbb{R}^{n, n}$.
    \end{proof}

\subsection{Prodotto matrice-matrice} \label{subsec:matmul}
    Date due matrici $A \in \mathbb{R}^{m, n}$ e $B \in \mathbb{R}^{n, k}$ il prodotto matrice-matrice (o prodotto matriciale) fra $A$ e $B$ è definito come:
    \[
        \begin{aligned}
        AB
        &=
        \begin{bmatrix}
            a_{11} & a_{12} & ... & a_{1n} \\
            a_{21} & a_{22} & ... & a_{2n} \\
            \vdots & \ddots & \ddots & \vdots \\
            a_{m1} & a_{m2} & ... & a_{mn}
        \end{bmatrix}
        \begin{bmatrix}
            b_{11} & b_{12} & ... & b_{1k} \\
            b_{21} & b_{22} & ... & b_{2k} \\
            \vdots & \ddots & \ddots & \vdots \\
            b_{n1} & b_{n2} & ... & b_{nk}
        \end{bmatrix}
        \\
        &=
        \begin{bmatrix}
            \displaystyle\sum_{i = 1}^{n}{a_{1i} b_{i1}} & \displaystyle\sum_{i = 1}^{n}{a_{1i} b_{i2}} & ... & \displaystyle\sum_{i = 1}^{n}{a_{1i} b_{ik}} \\
            \displaystyle\sum_{i = 1}^{n}{a_{2i} b_{i1}} & \displaystyle\sum_{i = 1}^{n}{a_{2i} b_{i2}} & ... & \displaystyle\sum_{i = 1}^{n}{a_{2i} b_{ik}} \\
            \vdots & \ddots & \ddots & \vdots \\
            \displaystyle\sum_{i = 1}^{n}{a_{mi} b_{i1}} & \displaystyle\sum_{i = 1}^{n}{a_{mi} b_{i2}} & ... & \displaystyle\sum_{i = 1}^{n}{a_{mi} b_{ik}}
        \end{bmatrix}
        \in \mathbb{R}^{m, k}
        \end{aligned}
    \]

    \vspace{10pt}
    Proprietà:
    \begin{enumerate}[label=(\arabic*)]
        \item Associativa: $\forall A \in \mathbb{R}^{m, n}, \forall B \in \mathbb{R}^{n, k}, \forall C \in \mathbb{R}^{k, p}, \; (AB) C = A (BC)$
        \item Associativa rispetto al prodotto scalare-matrice: $\forall \lambda \in \mathbb{R}, \forall A \in \mathbb{R}^{m, n}, \forall B \in \mathbb{R}^{n, k}, \; \lambda (A  B) = (\lambda A)  B = A  (\lambda B)$
        \item Distributiva sinistra rispetto alla somma fra matrici: $\forall A, B \in \mathbb{R}^{m, n}, \forall C \in \mathbb{R}^{k, m}, \; C (A + B) = C  A + C  B$
        \item Distributiva destra rispetto alla somma fra matrici: $\forall A, B \in \mathbb{R}^{m, n}, \forall C \in \mathbb{R}^{n, k}, \; (A + B) C = AC + BC$
        \item Distributiva della trasposizione rispetto al prodotto matriciale: $\forall A \in \mathbb{R}^{m, n}, \forall B \in \mathbb{R}^{n, k}, \; (A  B)^{T} = B^{T}  A^{T}$
    \end{enumerate}

\subsection{Matrici ridotte, pivot e rango di matrice} \label{subsec:matrank}
    Data una matrice $A \in \mathbb{R}^{m, n}$ essa si dice ridotta per righe (row echelon form) se e solo se il primo elemento non nullo
    di ogni riga, chiamato pivot, si trova a destra del pivot della riga superiore. $A$ si dice invece fortemente ridotta per righe,
    oppure in forma echelon (reduced row echelon form), se è ridotta per righe, ogni pivot vale 1 ed è l'unico elemento non nullo della colonna in cui si trova.

    \vspace{10pt}
    Un esempio di matrice ridotta per righe è:
    \[
        A =
        \begin{bmatrix}
            3 & 2 & 4 & 9 \\
            0 & 1 & 8 & 8 \\
            0 & 0 & 0 & 6 \\
            0 & 0 & 0 & 0
        \end{bmatrix}
    \]
    Mentre un esempio di matrice fortemente ridotta per righe:
    \[
        A =
        \begin{bmatrix}
            1 & 0 & 0 & 9 \\
            0 & 1 & 0 & 8 \\
            0 & 0 & 1 & 6 \\
            0 & 0 & 0 & 0
        \end{bmatrix}
    \]

    \vspace{10pt}
    Ogni matrice $A \in \mathbb{R}^{m, n}$ può essere portata in forma ridotta per righe attraverso un processo iterativo chiamato
    eliminazione di Gauss-Jordan, il quale consiste nell'applicare una delle seguenti traformazioni alle righe laddove è necessario:
    \begin{enumerate}[label=(E\textsubscript{\arabic*}), leftmargin=1cm]
        \item Sostituire una qualsiasi riga con la somma della riga stessa e un multiplo di un'altra riga. In simboli $\mathbf{R}_{i} + \lambda \mathbf{R}_{j} \rightarrow \mathbf{R}_{i}$ con $\lambda \in \mathbb{R}$
        \item Sostituire una qualsiasi riga con un multiplo non nullo della riga stessa. In simboli $\lambda \mathbf{R}_{i} \rightarrow \mathbf{R}_{i}$ con $\lambda \in \mathbb{R} - \{0\}$
        \item Sostituire una qualsiasi riga con un'altra. In simboli $\mathbf{R}_{j} \rightarrow \mathbf{R}_{i}$
    \end{enumerate}
    Dove $\mathbf{R}_{i}, \mathbf{R}_{j} \in \mathbb{R}^{1, n}$ sono le righe i e j-esime della matrice $A$.

    \vspace{10pt}
    Un esempio di eliminazione di Gauss è:
    \[
        A =
        \begin{bmatrix}
            3 & 1 & 4 & 9 \\
            3 & 2 & 8 & 8 \\
            0 & 3 & 0 & 6 \\
            0 & 0 & 0 & 7
        \end{bmatrix}
        \xrightarrow{\mathbf{R}_{2} - \mathbf{R}_{1} \rightarrow \mathbf{R}_{2} \; (E_1)}
        \begin{bmatrix}
            3 & 1 & 4 & 9 \\
            0 & 1 & 4 & -1 \\
            0 & 3 & 0 & 6 \\
            0 & 0 & 0 & 7
        \end{bmatrix}
        \xrightarrow{\mathbf{R}_{3} - 3 \mathbf{R}_{2} \rightarrow \mathbf{R}_{3} \; (E_1)}
        \begin{bmatrix}
            3 & 1 & 4 & 9 \\
            0 & 1 & 4 & -1 \\
            0 & 0 & -12 & 9 \\
            0 & 0 & 0 & 7
        \end{bmatrix}
    \]

    \vspace{10pt}
    \subsubsection{Definizione di rango di una matrice} \label{subsubsec:matrankdef}
    Il rango di una matrice è definito come il numero di righe (o colonne) linearmente indipendenti oppure, detto in altro modo, il rango di una matrice è la dimensione dello spazio generato dalle righe (o colonne) della matrice.

    \vspace{10pt}
    \begin{theorem}[Teorema del rango]
        Sia $A \in \mathbb{R}^{m, n}$ matrice e sia $R_A = \textnormal{span}(\mathbf{R}_{1}, \mathbf{R}_{2}, ..., \mathbf{R}_{m})$ lo spazio generato dalle righe di $A$, mentre $C_A = \textnormal{span}(\mathbf{C}_{1}, \mathbf{C}_{2}, ..., \mathbf{C}_{n})$ lo spazio
        generato dalle colonne di $A$, allora $\dim(R_A) = \dim(C_A)$
    \end{theorem}

    \begin{proof}[Dimostrazione]
        Sia $B_R = \{\mathbf{r}_{1}, ..., \mathbf{r}_{k}\}$ una base di $R_A$ con $\mathbf{r}_{i} \in B_R, \; \mathbf{r}_{i} = (r_{1}^{i}, ..., r_{n}^{i})$. E' allora possibile esprimere ogni riga come $\mathbf{R}_{i} = a_{1}^{i} \mathbf{r}_{1} + a_{2}^{i} \mathbf{r}_{2} + ... + a_{k}^{i} \mathbf{r}_{k}$.
        Cercando di visualizzare $A$ scomponendo ogni riga secondo $B_R$
        \[
            A =
            \begin{bmatrix}
                \mathbf{R}_{1} \\
                \mathbf{R}_{2} \\
                \vdots \\
                \mathbf{R}_{m}
            \end{bmatrix}
            =
            \begin{bmatrix}
                a_{1}^{1} r_{1}^{1} + a_{2}^{1} r_{1}^{2} + ... + a_{k}^{1} r_{1}^{k} & ... & a_{1}^{1} r_{n}^{1} + a_{2}^{1} r_{n}^{2} + ... + a_{k}^{1} r_{n}^{k} \\
                a_{1}^{2} r_{1}^{1} + a_{2}^{2} r_{1}^{2} + ... + a_{k}^{2} r_{1}^{k} & ... & a_{1}^{2} r_{n}^{1} + a_{2}^{2} r_{n}^{2} + ... + a_{k}^{2} r_{n}^{k} \\
                \vdots & \ddots & \vdots \\
                a_{1}^{m} r_{1}^{1} + a_{2}^{m} r_{1}^{2} + ... + a_{k}^{m} r_{1}^{k} & ... & a_{1}^{m} r_{n}^{1} + a_{2}^{m} r_{n}^{2} + ... + a_{k}^{m} r_{n}^{k}
            \end{bmatrix}
        \]
        In altre parole, qualisasi elemento $e_{ij}$ della matrice è nella forma $a_{1}^{i} r_{j}^{1} + a_{2}^{i} r_{j}^{2} + ... + a_{k}^{i} r_{j}^{k}$ dove il pedice di ogni termine $r$ rappresenta la colonna, e l'apice di ogni termine $a$ rappresenta la riga.

        Dunque, data una generica colonna $\mathbf{C}_{i}$
        \[
            \mathbf{C}_{i} =
            \begin{bmatrix}
                a_{1}^{1} \\
                a_{1}^{2} \\
                \vdots \\
                a_{1}^{m}
            \end{bmatrix}
            r_{i}^{1}
            +
            \begin{bmatrix}
                a_{2}^{1} \\
                a_{2}^{2} \\
                \vdots \\
                a_{2}^{m}
            \end{bmatrix}
            r_{i}^{2}
            +
            \dots
            +
            \begin{bmatrix}
                a_{k}^{1} \\
                a_{k}^{2} \\
                \vdots \\
                a_{k}^{m}
            \end{bmatrix}
            r_{i}^{k}
        \]

        il che implica che
        $G_C = \{
            \begin{bmatrix}
                a_{1}^{1} \\
                a_{1}^{2} \\
                \vdots \\
                a_{1}^{m}
            \end{bmatrix},
            \begin{bmatrix}
                a_{2}^{1} \\
                a_{2}^{2} \\
                \vdots \\
                a_{2}^{m}
            \end{bmatrix},
            ...,
            \begin{bmatrix}
                a_{k}^{1} \\
                a_{k}^{2} \\
                \vdots \\
                a_{k}^{m}
            \end{bmatrix}
            \}$
        sia un generatore di $C_A$. Quindi, per il lemma di

        \vspace{5pt}
        Steinitz, si ha che $\dim(C_A) \leq |G_C| = |B_R| = \dim(R_A)$. Si prenda in considerazione $A^T$ e si ripeta il procedimento svolto.
        Si giungerà alla conclusione che $\dim(C_{A^T}) = \dim(R_A) \leq \dim(R_{A^T}) = \dim(C_A)$, per cui $\dim(R_A) = \dim(C_A)$.
    \end{proof}

\vspace{10pt}
\subsection{Sistemi lineari} \label{subsec:matlinear}
    Un sistema lineare è un sistema di $m$ equazioni lineari a $n$ incognite della forma:
    \[
        \begin{cases}
            a_{11}x_1 + a_{12}x_2 + ... + a_{1n}x_n = b_1 \\
            a_{21}x_1 + a_{22}x_2 + ... + a_{2n}x_n = b_2 \\
            \vdots \\
            a_{m1}x_1 + a_{m2}x_2 + ... + a_{mn}x_n = b_m
        \end{cases}
    \]
    Per ogni sistema lineare è possibile costruire una sua versione "in matrice", ovvero:
    \[
        \begin{cases}
            a_{11}x_1 + a_{12}x_2 + ... + a_{1n}x_n = b_1 \\
            a_{21}x_1 + a_{22}x_2 + ... + a_{2n}x_n = b_2 \\
            \vdots \\
            a_{m1}x_1 + a_{m2}x_2 + ... + a_{mn}x_n = b_m
        \end{cases}
        \,
        \rightarrow
        \,
        \begin{bmatrix}
            a_{11} & a_{12} & ... & a_{1n} \\
            a_{21} & a_{22} & ... & a_{2n} \\
            \vdots & \ddots & \ddots & \vdots \\
            a_{m1} & a_{m2} & ... & a_{mn}
        \end{bmatrix}
        \,
        \begin{bmatrix}
            x_{1} \\
            x_{2} \\
            \vdots \\
            x_{n} \\
        \end{bmatrix}
        \,
        =
        \begin{bmatrix}
            b_{1} \\
            b_{2} \\
            \vdots \\
            b_{m} \\
        \end{bmatrix}
        \,
        \rightarrow
        A\mathbf{x} = \mathbf{b}
    \]
    La matrice aumentata del sistema $A | \mathbf{b} \in \mathbb{R}^{m, n+1}$ è definita come:
    \[
        \left[
        \begin{array}{cccc | c}
            a_{11} & a_{12} & ... & a_{1n} & b_1 \\
            a_{21} & a_{22} & ... & a_{2n} & b_2 \\
            \vdots & \ddots & \ddots & \vdots & \vdots \\
            a_{m1} & a_{m2} & ... & a_{mn} & b_m
        \end{array}
        \right]
    \]

    Un sistema lineare cui $\mathbf{b} = \mathbf{0}$ viene detto omogeneo ed è sempre risolubile in quanto ammette la soluzione triviale $\mathbf{x} = \mathbf{0}$.

    \vspace{10pt}
    \begin{theorem}[Teorema di Rouché-Capelli]
        Sia $S$ un sistema lineare e sia $A | \mathbf{b}$ la sua matrice aumentata. $S$ ha soluzione se e solo se
        $\textnormal{rank}(A | \mathbf{b})$ = $\textnormal{rank}(A)$ e la dimensione dello spazio delle soluzioni di $S$ vale $n - p$,
        con $n$ il numero di incognite del sistema e $p = \textnormal{rank}(A) = \textnormal{rank}(A | \mathbf{b})$
    \end{theorem}

    \begin{proof}[Dimostrazione]
        Il sistema $S$ può essere scritto in forma: $A\mathbf{x} = \mathbf{b}$. Si riscriva la relazione come: $L_A(\mathbf{x}) = \mathbf{b}$ dove $L_A: \mathbb{R}^n \rightarrow \mathbb{R}^m$ è un'applicazione lineare.
        $S$ ammette soluzione se e solo se $\mathbf{b} \in \textnormal{Im}(L_A)$, ovvero $\exists \mathbf{x} \in \mathbb{R}^n \; | \; L_A(\mathbf{x}) = \mathbf{b}$. E' possibile osservare come l'immagine di $L_A(\mathbf{x})$
        sia data da una combinazione lineare dei vettori colonna di $A$:
        \[
            A\mathbf{x} =
            \begin{bmatrix}
                a_{11}x_{1} + a_{12}x_{2} + \cdots + a_{1n}x_{n} \\
                a_{21}x_{1} + a_{22}x_{2} + \cdots + a_{2n}x_{n} \\
                \vdots \\
                a_{m1}x_{1} + a_{m2}x_{2} + \cdots + a_{mn}x_{n} \\
            \end{bmatrix}
            =
            \begin{bmatrix}
                a_{11} \\
                a_{21} \\
                \vdots \\
                a_{m1} \\
            \end{bmatrix}
            x_{1}
            +
            \begin{bmatrix}
                a_{12} \\
                a_{22} \\
                \vdots \\
                a_{m2} \\
            \end{bmatrix}
            x_{2}
            +
            \cdots
            +
            \begin{bmatrix}
                a_{1n} \\
                a_{2n} \\
                \vdots \\
                a_{mn} \\
            \end{bmatrix}
            x_{n}
        \]

        \vspace{10pt}
        Il che significa che $\mathbf{b}$ è contenuto nell'immagine di $L_A$ se e solo se $\mathbf{b}$ è nello span delle colonne di $A$.
        Si confronti il rango della matrice aumentata del sistema con il rango della matrice dei coefficienti. Dato che $A | \mathbf{b}$ è un'estensione di $A$, vi sono due casi:
        \begin{itemize}
            \item $\textnormal{rank}{(A | \mathbf{b})} = \textnormal{rank}(A)$
            \item $\textnormal{rank}{(A | \mathbf{b})} > \textnormal{rank}(A)$
        \end{itemize}
        Se $\textnormal{rank}{(A | \mathbf{b})} > \textnormal{rank}(A)$, allora $\mathbf{b}$ è linearmente indipendente ripetto a tutti i vettori di colonna di $A$, il che significa
        che $\mathbf{b} \notin \textnormal{Im}(L_A)$ e che $S$ non ha soluzione. Analogamente, se $\textnormal{rank}{(A | \mathbf{b})} = \textnormal{rank}(A)$, allora $\mathbf{b} \in \textnormal{Im}(L_A)$.

        \vspace{10pt}
        Appurato che esista almeno una soluzione $\mathbf{x}_{0}$, tutte le restanti soluzioni sono determinabili in maniera seguente:
        \[
            L_A(\mathbf{x}_{0}) = \mathbf{b} \implies \; \forall \mathbf{v} \in \ker(L_A), \; L_A(\mathbf{x}_{0} + \mathbf{v}) = L_A(\mathbf{x}_{0}) + L_A(\mathbf{v}) = \mathbf{b} + \mathbf{0} = \mathbf{b}
        \]
        Di conseguenza, la dimensione dello spazio di soluzione coincide con la dimensione del nucleo di $L_A$, la quale, per il teorema della dimensione, deve rispettare l'uguaglianza:
        \[
            \dim(\textnormal{Im}(L_A)) + \dim(\ker(L_A)) = \dim(\textnormal{Dom}(L_A)) \iff \textnormal{rank}(A) + \textnormal{null}(A) = n
        \]
        ricavando così la dimensione dello spazio di soluzione come $n - \textnormal{rank}(A)$.
    \end{proof}




\subsection{Equazioni matriciali} \label{subsec:matequation}
    Siano $A \in \mathbb{R}^{m, n}, \; X \in \mathbb{R}^{n, k}, \; B \in \mathbb{R}^{m, k}$, allora:
    \[
        AX = B
    \]
    è un'equazione matriciale la cui incognita è $X$.

    \vspace{10pt}
    Per le equazioni matriciali vige il teorema di Rouché-Capelli generalizzato, ovvero: \\
    Una qualsiasi eq. matriciale della forma $AX = B$ ha soluzione se e solo se rank$(A) = $ rank$(A | B)$ e qualora vi fossero delle soluzioni, la dimensione dello
    spazio di soluzione è data da $(n - p)k$ dove $p = \textnormal{rank}(A) = \textnormal{rank}(A | B)$.

    \vspace{10pt}
    \begin{corollary}[Corollario al teorema di Rouché-Capelli]
        Siano $A, X, B \in \mathbb{R}^{n, n}$ e sia $\textnormal{rank}(B) = n$, allora le seguenti condizioni sono equivalenti:
        \begin{enumerate}[label=(\Roman*)]
            \item $AX = B$ è risolubile
            \item $\textnormal{rank}(A) = n$
            \item La soluzione $X$ è unica
        \end{enumerate}
    \end{corollary}


\vspace{10pt}
\subsection{Matrici inverse, simili, ortogonali e special-ortogonali} \label{subsec:matspecial}
    \subsubsection{Matrici inverse} \label{subsubsec:matinv}
    Una matrice $A \in \mathbb{R}^{n, n}$ si dice invertibile se e solo se $\exists A^{-1} \in \mathbb{R}^{n, n} \; | \; AA^{-1} = I_n$. Se esiste la matrice inversa, allora essa è sia inversa a sinistra che inversa a destra, secondo quanto segue:
    \[
        C = A^{-1} \; (\textnormal{sx}), \; B = A^{-1} \; (\textnormal{dx}) \implies CA = AB = I \implies C = CI = CAB = IB = B
    \]
    Dato che una matrice quadrata può essere vista come un endomorfismo da $\mathbb{R}^{n}$ a $\mathbb{R}^{n}$ si può dire che esiste la matrice inversa di $A$ se e solo se l'endomorfismo è biettivo, quindi $\textnormal{rank}(A) = n$.

    \subsubsection{Matrici simili} \label{subsubsec:matalike}
    \vspace{10pt}
    Due matrici $A, B \in \mathbb{R}^{n, n}$ si dicono simili se:
    \[
        \exists P \in \mathbb{R}^{n, n} \; | \; A = P^{-1}BP, \;\: B = PAP^{-1}
    \]

    \subsubsection{Matrici ortogonali} \label{subsubsec:matorthogonal}
    Una matrice $A \in \mathbb{R}^{n, n}$ si dice ortogonale ($A \in O(n) \subset GL_{n}(\mathbb{R})$) se
    \[
        AA^{T} = A^{T}A = I
    \]
    $A$ è special-ortogonale ($A \in SO(n)$) se è ortogonale e $\det(A) = 1$.
    Il gruppo special-ortogonale $SO(n)$ può essere visto come il gruppo di tutte le rotazioni nello spazio euclideo $\mathbb{R}^{n}$.
